% =====================================================================
%  NX Cred — Tài liệu đặc tả ứng dụng
%  Biên dịch bằng XeLaTeX (bắt buộc, vì tiếng Việt có dấu):
%      xelatex nxcred-spec.tex
%      xelatex nxcred-spec.tex      (chạy hai lần để mục lục và tham chiếu đúng)
%  Hoặc tải lên Overleaf rồi chọn Compiler = XeLaTeX.
% =====================================================================
\documentclass[a4paper,11pt,oneside]{report}

\usepackage{fontspec}
\setmainfont{TeX Gyre Termes}
\setsansfont{TeX Gyre Heros}
\setmonofont{TeX Gyre Cursor}

\usepackage[margin=2.5cm]{geometry}
\usepackage{amsmath,amssymb,amsthm}
\usepackage{tikz}
\usetikzlibrary{arrows.meta,positioning,calc,fit,backgrounds}
\usepackage{enumitem}
\usepackage{booktabs}
\usepackage{longtable}
\usepackage{array}
\usepackage{xcolor}
\usepackage{titlesec}
\usepackage[hidelinks]{hyperref}

\definecolor{ink}{HTML}{16171A}
\definecolor{muted}{HTML}{6E7079}
\definecolor{accent}{HTML}{2F5FE0}
\definecolor{good}{HTML}{17795E}
\definecolor{soft}{HTML}{F5F5F1}
\definecolor{line}{HTML}{D8D8D2}

\hypersetup{
  colorlinks=true, linkcolor=accent, urlcolor=accent, citecolor=accent,
  pdftitle={NX Cred — Tài liệu đặc tả ứng dụng},
  pdfsubject={Ứng dụng định danh dùng zero-knowledge proof theo tiêu chuẩn AnonCreds}
}

\titleformat{\chapter}[display]
  {\sffamily\huge\bfseries\color{ink}}{\chaptertitlename\ \thechapter}{10pt}{\Huge}
\titlespacing*{\chapter}{0pt}{0pt}{26pt}
\titleformat{\section}{\sffamily\Large\bfseries\color{ink}}{\thesection}{0.6em}{}
\titleformat{\subsection}{\sffamily\large\bfseries\color{ink}}{\thesubsection}{0.6em}{}
\titleformat{\subsubsection}{\sffamily\normalsize\bfseries\color{ink}}{\thesubsubsection}{0.6em}{}

\theoremstyle{definition}
\newtheorem{definition}{Định nghĩa}[chapter]
\theoremstyle{plain}
\newtheorem{theorem}[definition]{Định lý}
\newtheorem{lemma}[definition]{Bổ đề}
\theoremstyle{remark}
\newtheorem*{remark}{Nhận xét}
\renewcommand{\proofname}{Chứng minh}

% ---- Cơ chế thuật ngữ -------------------------------------------------
% \gd{khoa}{Tên hiển thị}  : định nghĩa một mục trong chương Topology
% \g{khoa}{chữ hiển thị}   : mọi lần nhắc tới thuật ngữ, thành liên kết về Topology
\newcommand{\gd}[2]{\subsection{#2}\label{term:#1}}
\newcommand{\g}[2]{\hyperref[term:#1]{#2}}

\newcommand{\Zn}{\mathbb{Z}_n^{*}}
\newcommand{\Hid}{\mathcal{H}}
\newcommand{\Rev}{\mathcal{R}}
\newcommand{\Hash}{\mathsf{H}}
\newcommand{\enc}{\mathsf{enc}}

\tikzset{
  blk/.style={draw=line,fill=white,rounded corners=4pt,inner sep=6pt,
              font=\sffamily\scriptsize,align=center,minimum height=9mm},
  actor/.style={blk,fill=soft,font=\sffamily\small\bfseries,minimum width=34mm},
  chain/.style={blk,fill=soft,draw=accent,font=\sffamily\small\bfseries},
  step/.style={-{Stealth[length=2.2mm]},draw=ink,thick,font=\sffamily\tiny,align=center},
  dashx/.style={-{Stealth[length=2.2mm]},draw=muted,thick,dashed,font=\sffamily\tiny},
  note/.style={font=\sffamily\scriptsize\itshape,color=muted,align=center},
}

\begin{document}

% ======================= TRANG BÌA ====================================
\begin{titlepage}
\centering
\vspace*{3.2cm}
{\sffamily\Huge\bfseries NX Cred}\\[10pt]
{\sffamily\LARGE Tài liệu đặc tả ứng dụng}\\[26pt]
{\large\color{muted} Ứng dụng định danh sử dụng bằng chứng không tiết lộ\\
theo tiêu chuẩn AnonCreds}\\[46pt]
\rule{0.6\textwidth}{0.4pt}\\[16pt]
{\color{muted} Phiên bản demo \quad\textbullet\quad \today}
\vfill
\end{titlepage}

% ======================= ABSTRACT =====================================
\chapter*{Tóm tắt}
\addcontentsline{toc}{chapter}{Tóm tắt}

Đây là bộ tài liệu đặc tả về \textbf{NX Cred} --- ứng dụng định danh sử dụng
\g{zkp}{bằng chứng không tiết lộ (zero-knowledge proof)}, áp dụng tiêu chuẩn của
\g{anoncreds}{AnonCreds}.

Vấn đề mà NX Cred giải quyết: trong mô hình định danh hiện nay, mỗi lần một người chứng minh
danh tính là một lần dữ liệu cá nhân của họ bị sao chép thêm vào một cơ sở dữ liệu mới. Người
dùng mất quyền kiểm soát, bên tiếp nhận gánh rủi ro lưu trữ, và bên cấp giấy tờ vô tình trở
thành nơi nắm giữ toàn bộ lịch sử hoạt động của công dân.

NX Cred thay việc \emph{đưa dữ liệu} bằng việc \emph{đưa một bằng chứng toán học về dữ liệu}.
Hệ thống gồm ba bên --- \g{issuer}{Issuer}, \g{holder}{Holder} và \g{verifier}{Verifier} ---
cùng một \g{blockchain}{sổ cái công khai} giữ khoá công khai. Bên xác minh khẳng định được
tấm thẻ là thật, do đúng cơ quan có thẩm quyền ký, và đúng là của người đang đứng trước mặt
--- mà không nhìn thấy bất kỳ trường dữ liệu nào ngoài những trường người dùng chủ động mở.

Tài liệu mô tả đầy đủ: các khái niệm nền tảng, bài toán và lý do chọn giải pháp, bốn luồng
hoạt động giữa các bên, toàn bộ công thức của lõi mật mã, thiết kế bảo mật, thiết kế
blockchain, ba tình huống ứng dụng thực tế, cùng những giới hạn hiện tại và hướng phát triển.

\bigskip
\noindent\textbf{Cách đọc tài liệu này.} Mọi thuật ngữ chuyên môn, ở lần xuất hiện đầu tiên,
đều là một liên kết dẫn về chương~\ref{ch:topology} (\emph{Topology}) --- nơi thuật ngữ đó
được giải thích. Bạn có thể bấm vào bất kỳ từ nào chưa rõ để nhảy tới định nghĩa, rồi quay lại
đọc tiếp. Chương~\ref{ch:crypto} (lõi mật mã) có thể bỏ qua hoàn toàn nếu không cần đi sâu.

% ======================= MỤC LỤC ======================================
\tableofcontents

% ======================= TOPOLOGY =====================================
\chapter{Topology --- từ điển thuật ngữ}
\label{ch:topology}

Chương này là nơi tra cứu. Mỗi mục giải thích một thuật ngữ, và mọi chỗ khác trong tài liệu
nhắc tới thuật ngữ đó đều liên kết về đây.

\section{Khái niệm nền tảng}

\gd{ssi}{Self-sovereign identity (định danh tự chủ)}
Mô hình định danh trong đó người dùng --- chứ không phải một tổ chức nào --- là bên nắm giữ và
kiểm soát dữ liệu định danh của chính mình. Người dùng quyết định chia sẻ gì, với ai, trong bao
lâu. NX Cred được xây theo mô hình này.

\gd{anoncreds}{AnonCreds}
Viết tắt của \emph{Anonymous Credentials}: bộ tiêu chuẩn mô tả cách cấp và trình bày giấy tờ số
sao cho bên xác minh kiểm chứng được tính hợp lệ mà không đọc được dữ liệu, và không thể ghép
hai lần trình bày lại với nhau. Tiêu chuẩn này quy định bốn thành phần chính:
\g{linksecret}{link secret}, \g{clsig}{chữ ký CL}, \g{selective}{tiết lộ chọn lọc}, và
\g{unlink}{tính không liên kết được}. AnonCreds đã được vương quốc Bhutan áp dụng cho hệ thống
định danh quốc gia \g{ndi}{Bhutan NDI} phục vụ khoảng 200.000 người dân.

\gd{zkp}{Zero-knowledge proof (bằng chứng không tiết lộ)}
Thuật toán mật mã cho phép chứng minh một mệnh đề là đúng mà không tiết lộ bất cứ điều gì ngoài
chính tính đúng đắn đó. Một bằng chứng như vậy phải thoả ba tính chất: \emph{đầy đủ} (điều đúng
thì luôn chứng minh được), \emph{chắc chắn} (điều sai thì không thể chứng minh, kể cả khi kẻ
gian có năng lực tính toán lớn), và \emph{không tiết lộ} (bên kiểm kết thúc phiên mà không biết
thêm gì). Đây là trái tim của NX Cred.

\gd{ndi}{Bhutan NDI}
Hệ thống định danh số quốc gia của vương quốc Bhutan, xây trên tiêu chuẩn
\g{anoncreds}{AnonCreds}. Đây là minh chứng thực tế rằng mô hình này vận hành được ở quy mô
quốc gia chứ không chỉ trong phòng thí nghiệm.

\section{Các bên tham gia}

\gd{issuer}{Issuer --- bên cấp}
Bên có thẩm quyền cao nhất trong việc xác nhận danh tính: cơ quan quản lý căn cước, trường đại
học, hoặc tổ chức phát hành giấy tờ. Issuer đối chiếu hồ sơ \g{ekyc}{eKYC}, ký
\g{credential}{credential} bằng \g{blindsign}{chữ ký mù}, và công bố
\g{pubkey}{khoá công khai} cùng \g{schema}{schema} lên \g{blockchain}{blockchain}.

\gd{holder}{Holder --- ví người dùng}
Phần mềm phía người dùng, còn gọi là \emph{wallet}. Holder sinh và giữ
\g{linksecret}{link secret}, nhận và lưu \g{credential}{credential}, rồi tạo
\g{zkp}{bằng chứng không tiết lộ} mỗi khi cần chứng minh danh tính. Đây là bên duy nhất có thể
tạo ra bằng chứng hợp lệ.

\gd{verifier}{Verifier --- bên xác minh}
Bất kỳ tổ chức nào cần kiểm tra danh tính: ngân hàng, khách sạn, nhà tuyển dụng, hội đồng thi.
Verifier soạn \g{presreq}{yêu cầu trình bày} dựa trên \g{schema}{schema} lấy từ blockchain, rồi
kiểm bằng chứng nhận được bằng \g{pubkey}{khoá công khai} cũng lấy từ blockchain.

\gd{blockchain}{Blockchain --- sổ cái công khai}
Nơi \g{issuer}{issuer} công bố \g{pubkey}{khoá công khai} và \g{schema}{schema}. Điểm mấu chốt
là dữ liệu trên đó \emph{không ai sửa được}, kể cả chính issuer. Blockchain trong NX Cred
\textbf{không} lưu dữ liệu cá nhân, \textbf{không} lưu credential, và \textbf{không} ghi lại
các lần xác minh.

\section{Đại lượng mật mã}

\gd{bigint}{Big integer (số nguyên lớn)}
Số nguyên có độ dài hàng trăm đến hàng nghìn chữ số nhị phân --- lớn hơn rất nhiều so với kiểu
số thông thường của máy tính (thường 64 bit). Toàn bộ phép toán trong NX Cred làm việc trên
những số như vậy: \g{modulus}{modulus} dài khoảng 2048 bit, tức xấp xỉ 617 chữ số thập phân.
Độ lớn này chính là thứ khiến việc phá giải bằng cách thử mọi khả năng trở nên bất khả thi.

\gd{linksecret}{Link secret}
Một \g{bigint}{số nguyên lớn} ngẫu nhiên dài 256 bit, được sinh ra ngay sau khi người dùng tạo
ví, và \emph{không bao giờ rời khỏi thiết bị} --- kể cả \g{issuer}{issuer} cũng không nhìn
thấy. \g{credential}{Credential} được ký sao cho chỉ dùng được cùng với con số này. Hệ quả:
kẻ nào chép trộm được file credential mà không có link secret thì cũng không tạo được bằng
chứng hợp lệ. Mỗi tài khoản tương ứng với đúng một link secret duy nhất.

\gd{blinding}{Blinding factor (hệ số làm mù)}
Một \g{bigint}{số nguyên lớn} ngẫu nhiên do ví sinh ra để che giấu giá trị thật trong một
\g{commitment}{cam kết}. Vì mỗi giá trị bí mật giả định đều tương thích với đúng một blinding
factor, cam kết gửi đi không mang thông tin nào về giá trị thật. Ví giữ blinding factor lại để
về sau \g{unblind}{gỡ mù} chữ ký.

\gd{commitment}{Commitment (cam kết)}
Một số được tính từ giá trị bí mật cộng với \g{blinding}{blinding factor}, đóng vai như chiếc
phong bì dán kín: người khác ký được lên phong bì mà không đọc được nội dung bên trong, và
người gửi về sau không thể tráo đổi nội dung đó.

\gd{nonce}{Nonce (số dùng một lần)}
Số ngẫu nhiên do bên kiểm phát ra, chỉ dùng đúng một lần rồi huỷ. Nonce tham gia vào phép tính
\g{challenge}{thách thức}, nên một bằng chứng cũ gắn với nonce cũ không thể được phát lại cho
phiên mới. Đây là cơ chế chống tấn công phát lại (replay attack).

\gd{challenge}{Thách thức (challenge) và Fiat--Shamir}
Trong một giao thức chứng minh, bên kiểm gửi một câu hỏi ngẫu nhiên và bên chứng minh phải trả
lời. Biến đổi \emph{Fiat--Shamir} thay câu hỏi đó bằng một giá trị băm của chính dữ liệu bên
chứng minh vừa tạo ra. Vì không đoán trước được giá trị băm, bên chứng minh không thể dựng sẵn
câu trả lời giả --- nhờ vậy giao thức từ hai chiều trở thành một chiều, gửi một lần là xong.

\gd{sigma}{Sigma protocol (giao thức sigma)}
Dạng giao thức ba bước --- cam kết, thách thức, đáp số --- dùng để chứng minh \emph{biết} một
bí mật mà không nói ra bí mật đó. NX Cred dùng sigma protocol ở bước xin cấp credential, để ví
chứng minh mình thật sự nắm \g{linksecret}{link secret} và \g{blinding}{blinding factor} nằm
bên trong \g{commitment}{cam kết} vừa gửi.

\gd{clsig}{CL signature (chữ ký Camenisch--Lysyanskaya)}
Lược đồ chữ ký số cho phép ký lên nhiều thuộc tính cùng lúc, ký lên cả phần bị che, và cho phép
người giữ chữ ký về sau \g{rerand}{ngẫu nhiên hoá} lại nó mà vẫn giữ tính hợp lệ. Chính hai
tính chất cuối khiến CL signature phù hợp cho định danh ẩn danh, điều mà chữ ký RSA hay ECDSA
thông thường không làm được.

\gd{blindsign}{Blind signing (ký mù)}
Cách ký trong đó bên ký không nhìn thấy một phần nội dung mình đang ký, nhưng chữ ký vẫn ràng
buộc chặt vào phần bị che đó. Trong NX Cred, \g{issuer}{issuer} ký lên
\g{commitment}{cam kết} chứa \g{linksecret}{link secret} mà không hề biết link secret là gì.

\gd{unblind}{Unblinding (gỡ mù)}
Thao tác ví thực hiện sau khi nhận chữ ký từ issuer: dùng \g{blinding}{blinding factor} đã giữ
lại để biến chữ ký trên cam kết thành chữ ký hoàn chỉnh trên dữ liệu thật.

\gd{credential}{Credential (thẻ định danh)}
Kết quả cuối cùng của quá trình cấp: bộ ba số $(A, e, v)$ tạo thành chữ ký của issuer lên toàn
bộ thuộc tính cộng với \g{linksecret}{link secret}. Đây là thứ ví lưu lại và dùng để tạo bằng
chứng về sau.

\gd{schema}{Schema (khuôn mẫu giấy tờ)}
Danh sách các trường mà một loại giấy tờ có. Ví dụ với thẻ sinh viên: mã số sinh viên, họ và
tên, mã lớp, ngày hết hạn. Schema được \g{issuer}{issuer} công bố lên
\g{blockchain}{blockchain}, và nó là ranh giới cứng: issuer không ký gì ngoài các trường trong
schema, \g{verifier}{verifier} cũng không hỏi được gì ngoài các trường đó.

\gd{pubkey}{Khoá công khai (credential definition)}
Bộ số issuer công bố để người khác kiểm chữ ký của mình. Trong NX Cred gồm:
\g{modulus}{modulus} $n$, ba cơ số $S, R, Z$, và một cơ số riêng $R_i$ cho mỗi trường trong
\g{schema}{schema}. Chi tiết xem mục~\ref{sec:setup}.

\gd{modulus}{Modulus}
Số $n = p \cdot q$ với $p, q$ là hai \g{safeprime}{safe prime} bí mật. Mọi phép nhân và luỹ
thừa trong hệ thống đều thực hiện modulo $n$. Biết $n$ nhưng không biết $p, q$ thì không ký
được --- đó là nền tảng an toàn của lược đồ.

\gd{safeprime}{Safe prime}
Số nguyên tố dạng $p = 2p' + 1$ trong đó $p'$ cũng là số nguyên tố. Dùng safe prime khiến nhóm
tính toán có bậc chứa thừa số nguyên tố lớn, chặn được các thuật toán tấn công khai thác nhóm
có bậc phân tích được thành nhiều thừa số nhỏ.

\gd{rerand}{Randomization (ngẫu nhiên hoá chữ ký)}
Thao tác ví thực hiện trước mỗi lần trình bày: biến chữ ký $A$ thành $A' = A \cdot S^{r}$ với
$r$ ngẫu nhiên mới. Chữ ký gửi đi vì thế mỗi lần một khác, tạo ra
\g{unlink}{tính không liên kết được}.

\section{Tính chất và cơ chế}

\gd{selective}{Selective disclosure (tiết lộ chọn lọc)}
Khả năng mở đúng những trường được hỏi trên một tấm thẻ và giữ kín phần còn lại, trong khi bên
xác minh vẫn kiểm chứng được tấm thẻ là trọn vẹn và hợp lệ. Các trường không mở vẫn tham gia
vào phép tính dưới dạng đã che.

\gd{unlink}{Unlinkability (tính không liên kết được)}
Tính chất khiến hai lần trình bày cùng một credential không thể bị ghép lại với nhau. Đạt được
nhờ \g{rerand}{ngẫu nhiên hoá} chữ ký trước mỗi lần gửi: không có mã số cố định nào đi theo
người dùng, nên hai verifier gộp dữ liệu cũng không tìm ra điểm chung.

\gd{presreq}{Presentation request (yêu cầu trình bày)}
Bản yêu cầu verifier gửi cho ví, gồm: những trường cần được tiết lộ, những mệnh đề cần được
chứng minh mà không tiết lộ, và một \g{nonce}{nonce}.

\gd{predicate}{Predicate proof (bằng chứng vị từ)}
Dạng bằng chứng chứng minh một quan hệ so sánh trên thuộc tính số mà không lộ giá trị --- ví dụ
``năm sinh $\le$ 2007'' để chứng minh đủ 18 tuổi. NX Cred hiện \emph{chưa} hiện thực phần này;
xem mục~\ref{sec:limits}.

\gd{ekyc}{eKYC}
Viết tắt của \emph{electronic Know Your Customer}: bước đối chiếu danh tính điện tử, trong đó
người dùng quét giấy tờ và bên cấp so với hồ sơ mình đang giữ. Trong NX Cred, đây là bước duy
nhất dữ liệu định danh được đối chiếu trực tiếp, và nó chỉ diễn ra \emph{một lần} trong đời một
tài khoản.

\gd{passkey}{Passkey}
Khoá sinh trắc học theo chuẩn WebAuthn (vân tay, khuôn mặt, mã PIN thiết bị) dùng để mở kho lưu
trữ của ví. Cần phân biệt rõ: passkey là lớp bảo vệ \emph{quyền truy cập} vào thiết bị, không
phải cơ chế mật mã tạo nên bằng chứng.

\gd{revocation}{Revocation (thu hồi)}
Cơ chế cho phép issuer vô hiệu hoá một credential đã cấp. NX Cred hiện \emph{chưa} có phần này;
xem mục~\ref{sec:limits}.

% ======================= INTRODUCTION =================================
\chapter{Giới thiệu}

\section{AnonCreds là tiêu chuẩn gì}

\g{anoncreds}{AnonCreds} là bộ tiêu chuẩn mô tả cách cấp phát và trình bày giấy tờ số sao cho
ba yêu cầu sau đồng thời được thoả mãn:

\begin{enumerate}[leftmargin=1.5em,itemsep=5pt]
  \item \textbf{Bên xác minh tin được.} Bằng chứng nhận về phải chứng tỏ giấy tờ do đúng cơ
        quan có thẩm quyền ký, chưa bị sửa, và thuộc về đúng người đang trình.
  \item \textbf{Người dùng không lộ thừa.} Chỉ những trường được hỏi và được người dùng đồng ý
        mới hiện ra; phần còn lại vẫn nằm trong phép tính nhưng ở dạng đã che.
  \item \textbf{Không ai theo dõi được.} Hai lần trình cùng một tấm thẻ phải không ghép lại
        được với nhau, và bên cấp giấy tờ không được biết thẻ đang dùng ở đâu.
\end{enumerate}

Ba yêu cầu này xung đột nhau nếu dùng chữ ký số thông thường. Chữ ký RSA hay ECDSA là một giá
trị cố định: đưa ra lần nào cũng như lần nào, nên nó chính là một mã số theo dõi hoàn hảo. Muốn
giấu bớt thuộc tính thì chữ ký không còn kiểm được. AnonCreds giải quyết bằng bốn thành phần:

\begin{description}[leftmargin=0pt,style=nextline,itemsep=7pt]
  \item[\g{linksecret}{Link secret}] gắn credential với đúng một ví, khiến bản sao bị đánh cắp
    trở nên vô dụng.
  \item[\g{blindsign}{Ký mù}] cho phép issuer ký lên link secret mà không nhìn thấy nó.
  \item[\g{selective}{Tiết lộ chọn lọc}] cho phép mở từng trường riêng lẻ.
  \item[\g{unlink}{Ngẫu nhiên hoá}] khiến mỗi lần trình bày là một lần hoàn toàn mới.
\end{description}

Tiêu chuẩn này không chỉ nằm trên giấy. Vương quốc Bhutan đã triển khai
\g{ndi}{Bhutan NDI} trên nền AnonCreds, phục vụ khoảng 200.000 người dân.

\section{Zero-knowledge proof là thuật toán gì}

\g{zkp}{Zero-knowledge proof} là thuật toán cho phép chứng minh một mệnh đề đúng mà không tiết
lộ vì sao nó đúng.

Để hình dung, hãy xét một ví dụ không cần tới toán. Có một hang động hình vòng, lối vào chia
làm hai nhánh trái và phải, hai nhánh nối nhau ở cuối bằng một cánh cửa khoá. Bạn tuyên bố mình
biết mật khẩu cánh cửa đó. Người kiểm đứng ngoài, không thấy bạn chọn nhánh nào. Bạn đi vào một
nhánh tuỳ ý; sau đó người kiểm hô ``trái'' hoặc ``phải'', và bạn phải đi ra bằng nhánh được hô.

Nếu bạn thật sự biết mật khẩu, bạn luôn ra đúng nhánh. Nếu không biết, bạn chỉ ra đúng khi may
mắn đã chọn sẵn nhánh đó --- xác suất 50\%. Lặp lại 20 lần, xác suất ăn may là $2^{-20}$, tức
khoảng một phần triệu. Người kiểm buộc phải tin bạn biết mật khẩu, dù chưa từng nghe mật khẩu.

Ví dụ này chứa đủ ba tính chất của một bằng chứng không tiết lộ:

\begin{itemize}[leftmargin=1.5em,itemsep=4pt]
  \item \textbf{Đầy đủ} --- người biết mật khẩu luôn qua được.
  \item \textbf{Chắc chắn} --- người không biết gần như chắc chắn bị phát hiện.
  \item \textbf{Không tiết lộ} --- người kiểm không học được gì về mật khẩu.
\end{itemize}

Điểm mấu chốt là bên chứng minh phải phản ứng với một \g{challenge}{thách thức ngẫu nhiên} mà
họ không đoán trước được. Trong hệ thống thật, thách thức đó là một giá trị băm, và việc lặp
lại 20 lần được thay bằng việc dùng một con số đủ lớn --- 256 bit --- để xác suất ăn may nhỏ
tới mức không đáng kể.

\section{Giới thiệu bài toán}

\subsection{User story}

\textbf{Minh là sinh viên năm cuối.} Trong sáu tháng, Minh phải chứng minh danh tính ở tám nơi
khác nhau: mở tài khoản ngân hàng để nhận lương thực tập, ký hợp đồng thuê trọ, làm thẻ thư
viện liên trường, đăng ký thi chứng chỉ ngoại ngữ, nhận phòng khách sạn khi đi thực tập xa, mua
sim điện thoại, đăng ký dịch vụ ví điện tử, và nộp hồ sơ xin việc.

Mỗi lần như vậy, Minh đưa căn cước ra. Bên kia chụp lại hoặc photo. Sau sáu tháng, dữ liệu định
danh đầy đủ của Minh --- số căn cước, ngày sinh, quê quán, địa chỉ thường trú --- nằm trong tám
cơ sở dữ liệu mà Minh không kiểm soát, không biết được bảo vệ ra sao, và không thể yêu cầu xoá.

Ba vấn đề lộ ra:

\begin{description}[leftmargin=0pt,style=nextline,itemsep=7pt]
  \item[Nhân bản không thu hồi được]
    Tám bản sao nghĩa là tám cơ hội rò rỉ. Chỉ cần một nơi bị tấn công là dữ liệu của Minh ra
    ngoài, và không có cách nào ``thu hồi'' một tấm căn cước đã bị chụp. Tệ hơn, Minh thường
    không biết vụ rò rỉ đã xảy ra.

  \item[Lộ thừa thông tin]
    Khách sạn chỉ cần biết Minh là người thật và đủ tuổi thuê phòng. Nhưng tấm căn cước nói cho
    họ luôn cả quê quán và địa chỉ nhà. Minh trả lời một câu hỏi bằng cách tiết lộ mười câu trả
    lời --- và không có lựa chọn nào khác.

  \item[Giám sát phát sinh ngoài ý muốn]
    Ở những hệ thống có kiểm tra trực tuyến, mỗi lần Minh dùng giấy tờ là một lần bên tiếp nhận
    hỏi lại cơ quan cấp. Lặp lại đủ nhiều, cơ quan cấp trở thành nơi nắm được bản đồ hoạt động
    của Minh: khi nào mở tài khoản, thuê nhà ở đâu, đi những đâu. Không ai thiết kế hệ thống với
    mục đích giám sát, nhưng kết quả vẫn là giám sát.
\end{description}

\subsection{Vì sao NX Cred là giải pháp}

NX Cred đảo ngược mô hình: thay vì đưa dữ liệu rồi mong bên kia giữ gìn, người dùng đưa một
\g{zkp}{bằng chứng toán học} về dữ liệu đó.

\begin{center}
\begin{tabular}{@{}p{5.2cm}p{4.2cm}p{5.2cm}@{}}
\toprule
\textbf{\sffamily Vấn đề} & \textbf{\sffamily Cơ chế} & \textbf{\sffamily Kết quả} \\
\midrule
Nhân bản dữ liệu &
\g{zkp}{Bằng chứng} thay cho bản sao &
Bên tiếp nhận không có gì để làm rò rỉ \\
\addlinespace
Lộ thừa thông tin &
\g{selective}{Tiết lộ chọn lọc} &
Chỉ mở đúng trường được hỏi \\
\addlinespace
Giám sát phát sinh &
\g{pubkey}{Khoá công khai} trên \g{blockchain}{blockchain} &
Issuer không cần và không thể biết \\
\addlinespace
Thẻ bị đánh cắp &
\g{linksecret}{Link secret} trên máy &
Bản sao credential vô dụng \\
\addlinespace
Bị ghép dấu vết &
\g{rerand}{Ngẫu nhiên hoá} mỗi lần &
Hai lần trình không liên kết được \\
\bottomrule
\end{tabular}
\end{center}

Quay lại câu chuyện của Minh: sau khi dùng NX Cred, Minh làm \g{ekyc}{eKYC} đúng \emph{một
lần} với cơ quan cấp. Tám nơi còn lại nhận về tám bằng chứng toán học khác nhau. Không nơi nào
giữ dữ liệu gốc, không nơi nào ghép được với nơi khác, và cơ quan cấp không biết Minh đã đi
những đâu.

% ======================= LUỒNG HOẠT ĐỘNG ==============================
\chapter{Mô tả luồng hoạt động}
\label{ch:flows}

\section{Tổng quan luồng hoạt động}

Toàn hệ thống vận hành qua bảy bước, chia làm hai giai đoạn tách biệt hoàn toàn: \emph{cấp thẻ}
diễn ra một lần, \emph{xác minh} diễn ra mỗi khi cần.

\begin{enumerate}[leftmargin=1.6em,itemsep=6pt]
  \item Người dùng bấm \g{ekyc}{eKYC} trong ví để yêu cầu được cấp thẻ. Ví gửi dữ liệu giấy tờ
        cho \g{issuer}{issuer} đối chiếu.

  \item Issuer đối chiếu hồ sơ và phát về một \g{nonce}{nonce}, đồng thời yêu cầu ví tạo một
        \g{zkp}{bằng chứng} chứng minh nó thật sự sở hữu \g{linksecret}{link secret} và
        \g{blinding}{blinding factor} nằm trong \g{commitment}{cam kết} sắp gửi.

  \item Ví gửi lại nonce kèm bằng chứng đó (\g{sigma}{sigma protocol}). Issuer kiểm bằng chứng;
        nếu hợp lệ thì thực hiện \g{clsig}{CL signature}, tức \g{blindsign}{ký mù} lên cam kết
        và toàn bộ thuộc tính.

  \item Ví nhận chữ ký về, \g{unblind}{gỡ mù}, tự kiểm lại, rồi lưu thành
        \g{credential}{credential}. Giai đoạn cấp thẻ kết thúc tại đây và không lặp lại.

  \item Khi cần xác minh, \g{verifier}{verifier} tạo mã QR chứa
        \g{presreq}{yêu cầu trình bày} kèm một \g{nonce}{nonce} mới.

  \item Ví \g{rerand}{ngẫu nhiên hoá} chữ ký, làm mù các thuộc tính không chia sẻ, rồi tạo bằng
        chứng. \emph{Mỗi lần yêu cầu là một lần làm mù lại và tạo lại bằng chứng từ đầu.}

  \item Verifier nhận bằng chứng, khôi phục lại giá trị cần so sánh và đối chiếu. Song song đó,
        issuer đã công bố \g{pubkey}{khoá công khai} và \g{schema}{schema} lên
        \g{blockchain}{blockchain}; verifier lấy từ đó để biết cần hỏi gì và kiểm bằng cách nào.
\end{enumerate}

\begin{center}
\begin{tikzpicture}[node distance=15mm]
  \node[chain] (chain) {Blockchain\\[1pt]\tiny\mdseries khoá công khai + schema};
  \node[actor,below left=18mm and 24mm of chain] (issuer) {Issuer};
  \node[actor,below=32mm of chain] (holder) {Holder / Wallet};
  \node[actor,below right=18mm and 24mm of chain] (verifier) {Verifier};

  \draw[step] (issuer) -- node[left,xshift=-1mm]{công bố\\một lần} (chain);
  \draw[step] (chain) -- node[right,xshift=1mm]{đọc\\khi cần} (verifier);
  \draw[step] (issuer) -- node[below,sloped]{cấp thẻ (một lần)} (holder);
  \draw[step] (holder) -- node[below,sloped]{trình bằng chứng (mỗi lần)} (verifier);
  \draw[dashx] (issuer) to[out=-40,in=220] node[below,note,yshift=-1mm]{không có kênh trực tiếp} (verifier);
\end{tikzpicture}
\end{center}

\section{Luồng chi tiết: Issuer và Wallet}
\label{sec:flow-iw}

Đây là giai đoạn cấp thẻ, chạy đúng một lần cho mỗi người dùng.

\begin{center}
\begin{tikzpicture}[node distance=7mm]
  \node[actor] (h) {Wallet};
  \node[actor,right=58mm of h] (i) {Issuer};

  \node[blk,below=7mm of h,text width=40mm] (h1) {1. Người dùng bấm eKYC,\\ quét giấy tờ};
  \node[blk,below=7mm of i,text width=40mm] (i1) {2. Đối chiếu hồ sơ,\\ phát nonce};
  \node[blk,below=7mm of h1,text width=40mm] (h2) {3. Sinh link secret\\ và blinding factor};
  \node[blk,below=7mm of h2,text width=40mm] (h3) {4. Tính cam kết $u$\\ + bằng chứng sigma};
  \node[blk,below=19mm of i1,text width=40mm] (i2) {5. Kiểm bằng chứng};
  \node[blk,below=7mm of i2,text width=40mm] (i3) {6. CL signature:\\ ký mù $(A,e,v'')$};
  \node[blk,below=14mm of h3,text width=40mm] (h4) {7. Gỡ mù, tự kiểm,\\ lưu credential};

  \draw[step] (h1.east) -- node[above]{dữ liệu eKYC} (i1.west);
  \draw[step] (i1.south west) to[out=250,in=20] node[above,pos=.6]{nonce} (h2.east);
  \draw[step] (h3.east) -- node[above]{$(\text{nonce}, u, c, \hat v, \widehat{ls})$} (i2.west);
  \draw[step] (i2) -- (i3);
  \draw[step] (i3.south west) to[out=250,in=20] node[above,pos=.55]{$(A, e, v'')$} (h4.east);
\end{tikzpicture}
\end{center}

\subsection{Vì sao cần bằng chứng ở bước 4}

Nếu ví chỉ gửi \g{commitment}{cam kết} $u$ mà không kèm bằng chứng, một kẻ gian có thể chặn
bắt cam kết của người khác rồi nộp như của mình. Hắn sẽ nhận về một credential gắn với
\g{linksecret}{link secret} mà hắn không biết --- vô dụng với hắn, nhưng gây rối loạn sổ sách và
tiêu mất suất cấp thẻ của nạn nhân. \g{sigma}{Sigma protocol} ở bước này buộc người nộp phải
thật sự biết cặp số nằm bên trong cam kết.

\subsection{Vì sao ví phải tự kiểm ở bước 7}

Ví kiểm lại chữ ký trước khi lưu. Nếu issuer trả về một chữ ký hỏng --- do lỗi hoặc cố ý --- ví
phát hiện ngay tại chỗ, thay vì đến lúc đi xác minh mới hỏng và không biết quy trách nhiệm cho
ai.

\section{Luồng chi tiết: Wallet và Verifier}
\label{sec:flow-wv}

Đây là luồng chính, chạy xuyên suốt vòng đời ứng dụng.

\begin{center}
\begin{tikzpicture}[node distance=7mm]
  \node[actor] (v) {Verifier};
  \node[actor,right=58mm of v] (h) {Wallet};

  \node[blk,below=7mm of v,text width=40mm] (v1) {1. Soạn yêu cầu theo schema,\\ tạo QR + nonce};
  \node[blk,below=7mm of h,text width=40mm] (h1) {2. Người dùng quét QR,\\ xem yêu cầu};
  \node[blk,below=7mm of h1,text width=40mm] (h2) {3. Chọn trường chia sẻ};
  \node[blk,below=7mm of h2,text width=40mm] (h3) {4. Ngẫu nhiên hoá chữ ký\\ $A' = A \cdot S^{r}$};
  \node[blk,below=7mm of h3,text width=40mm] (h4) {5. Làm mù link secret\\ và thuộc tính ẩn};
  \node[blk,below=7mm of h4,text width=40mm] (h5) {6. Tạo bằng chứng};
  \node[blk,below=36mm of v1,text width=40mm] (v2) {7. Lấy khoá công khai\\ từ blockchain};
  \node[blk,below=7mm of v2,text width=40mm] (v3) {8. Khôi phục và so sánh};

  \draw[step] (v1.east) -- node[above]{QR + nonce} (h1.west);
  \draw[step] (h1) -- (h2);
  \draw[step] (h2) -- (h3);
  \draw[step] (h3) -- (h4);
  \draw[step] (h4) -- (h5);
  \draw[step] (h5.south west) to[out=250,in=20] node[above,pos=.55]{bằng chứng} (v2.east);
  \draw[step] (v2) -- (v3);
\end{tikzpicture}
\end{center}

Điểm quan trọng nhất của luồng này: các bước 4--6 chạy lại \emph{từ đầu} cho mỗi yêu cầu. Ví
không lưu sẵn một bằng chứng rồi dùng lại. Nhờ vậy hai verifier khác nhau nhận về hai bằng
chứng không có điểm chung nào, và một bằng chứng bị bắt trộm cũng không dùng lại được cho
\g{nonce}{nonce} khác.

Về mặt sản phẩm, bước 3 đáng nhấn mạnh: quyền quyết định thuộc về người dùng. Verifier chỉ được
\emph{hỏi}; ví hỏi lại chủ của nó trước khi trả lời. Những trường không nằm trong yêu cầu thì bị
khoá cứng trên giao diện --- người dùng không thể vô tình chia sẻ.

\section{Luồng chi tiết: Verifier và Issuer}
\label{sec:flow-vi}

Đây là mục ngắn nhất tài liệu, vì nội dung của nó là một sự vắng mặt: \textbf{verifier và issuer
không có kênh liên lạc trực tiếp nào.}

Trong mô hình truyền thống, verifier phải gọi về issuer để hỏi ``giấy tờ này có thật không''.
Câu hỏi đó tạo ra ba hệ quả xấu: issuer biết được toàn bộ hoạt động của người dùng; hệ thống
sập nếu issuer ngừng hoạt động; và mỗi lần xác minh phải chờ một vòng mạng.

NX Cred cắt đứt kênh đó bằng cách để issuer công bố \g{pubkey}{khoá công khai} lên
\g{blockchain}{blockchain} một lần duy nhất. Verifier tự đọc, tự kiểm, không cần xin phép và
không cần issuer còn trực tuyến. Issuer không biết --- và về mặt kỹ thuật là \emph{không thể}
biết --- thẻ đang được dùng ở đâu.

\section{Luồng chi tiết: Issuer, Blockchain và Verifier}
\label{sec:flow-ibv}

\begin{center}
\begin{tikzpicture}[node distance=18mm]
  \node[actor] (i) {Issuer};
  \node[chain,right=30mm of i,text width=34mm] (c) {Blockchain\\[2pt]\tiny\mdseries CredentialRegistry};
  \node[actor,right=30mm of c] (v) {Verifier};

  \draw[step] (i) -- node[above,align=center]{ghi một lần\\\tiny khoá công khai + schema} (c);
  \draw[step] (c) -- node[above,align=center]{đọc khi cần\\\tiny không cần xin phép} (v);
\end{tikzpicture}
\end{center}

Quy trình gồm hai thao tác, cách nhau về thời gian có thể là hàng tháng:

\begin{enumerate}[leftmargin=1.6em,itemsep=5pt]
  \item \textbf{Issuer ghi lên chain.} Khi thiết lập hệ thống, issuer sinh cặp khoá rồi ghi
        phần công khai cùng \g{schema}{schema} vào hợp đồng. Thao tác này thực hiện một lần.

  \item \textbf{Verifier đọc từ chain.} Khi cần xác minh, verifier đọc schema để biết loại giấy
        tờ này có những trường nào --- từ đó soạn được yêu cầu hợp lệ --- và đọc khoá công khai
        để kiểm bằng chứng.
\end{enumerate}

Điểm cốt lõi về mặt an ninh: verifier \textbf{không bao giờ} nhận khoá công khai từ bên đang
trình bằng chứng. Nếu nhận, kẻ gian chỉ việc tự sinh cặp khoá của mình, tự ký một tấm thẻ giả,
rồi nộp kèm khoá công khai của chính hắn --- mọi phép kiểm sẽ đều qua. Lấy khoá từ một nguồn
độc lập và không sửa được là thứ duy nhất chặn đường tấn công này.

% ======================= LÕI MẬT MÃ ===================================
\chapter{Lõi mật mã ứng dụng}
\label{ch:crypto}

\begin{quote}\itshape\color{muted}
Chương này có thể bỏ qua hoàn toàn nếu bạn không muốn tìm hiểu sâu. Toàn bộ phần còn lại của
tài liệu đọc được mà không cần tới đây.
\end{quote}

Mọi phép nhân và luỹ thừa dưới đây đều hiểu là trong $\Zn$, tức modulo
\g{modulus}{$n$}, trừ khi ghi rõ là phép toán trên $\mathbb{Z}$.

\section{Tham số và thiết lập khoá}
\label{sec:setup}

Issuer chọn hai \g{safeprime}{safe prime} $p, q$ dài 1024 bit, đặt $n = pq$ và
$\varphi = (p-1)(q-1)$. Lấy ngẫu nhiên $a, b, z, r_i \in [0, n)$ rồi đặt
\[
S = a^2, \qquad R = b^2, \qquad Z = z^2, \qquad R_i = r_i^2 \pmod n .
\]

Bình phương mọi cơ số đưa chúng vào nhóm các thặng dư bình phương, bảo đảm mọi phần tử tham gia
phép kiểm nằm cùng một nhóm con. Bộ $(n, S, R, Z, \{R_i\})$ chính là
\g{pubkey}{khoá công khai} được ghi lên \g{blockchain}{blockchain}; cặp $(p, q)$ giữ kín tuyệt
đối.

Thuộc tính được mã hoá thành số bằng hàm
\[
\enc(x) =
\begin{cases}
\operatorname{int}(x) & \text{nếu } x \text{ là chuỗi toàn chữ số},\\
\operatorname{int}(\mathrm{SHA\text{-}256}(x)) & \text{ngược lại,}
\end{cases}
\]
và \g{challenge}{thách thức} tính bằng
$\Hash(x, y, \mathit{nonce}) = \operatorname{int}(\mathrm{SHA\text{-}256}(x \| y \| \mathit{nonce}))$.

\section{Công thức giữa Issuer và Wallet}

\subsection{Cam kết}

Ví sinh \g{linksecret}{link secret} $\mathit{ls}$ dài 256 bit và
\g{blinding}{blinding factor} $v'$ dài 2048 bit, rồi tính
\begin{equation}
u = S^{v'} R^{\mathit{ls}} \bmod n .
\end{equation}

\begin{theorem}[Cam kết giấu kín hoàn hảo]
Với mỗi giá trị giả định $\widehat{\mathit{ls}}$ đều tồn tại đúng một $\widehat{v}$ sao cho
$S^{\widehat v} R^{\widehat{\mathit{ls}}} = u$. Do đó $u$ không mang thông tin nào về
$\mathit{ls}$.
\end{theorem}

\begin{proof}
Gọi $k$ là bậc của $S$ và giả sử $R = S^{\alpha}$. Khi đó $u = S^{v' + \alpha \mathit{ls}}$.
Với $\widehat{\mathit{ls}}$ bất kỳ, chọn
$\widehat v \equiv v' + \alpha\mathit{ls} - \alpha\widehat{\mathit{ls}} \pmod k$ thì
$S^{\widehat v} R^{\widehat{\mathit{ls}}} = S^{v' + \alpha\mathit{ls}} = u$. Nghiệm duy nhất
modulo $k$, nên phân phối của $u$ độc lập với $\mathit{ls}$.
\end{proof}

\subsection{Sigma protocol}

Ví lấy $\tilde v$ dài 3488 bit, $\widetilde{\mathit{ls}}$ dài 593 bit rồi tính
\[
\tilde u = S^{\tilde v} R^{\widetilde{\mathit{ls}}}, \qquad
c = \Hash(u, \tilde u, \mathit{nonce}), \qquad
\hat v = \tilde v + c v', \qquad
\widehat{\mathit{ls}} = \widetilde{\mathit{ls}} + c\,\mathit{ls},
\]
với hai phép cộng cuối thực hiện trên $\mathbb{Z}$. Ví gửi
$(\mathit{nonce}, u, c, \hat v, \widehat{\mathit{ls}})$ và \emph{không} gửi $\tilde u$.

Issuer kiểm $1 < u < n$, $\gcd(u,n) = 1$, rồi tính
$\tilde u' = S^{\hat v} R^{\widehat{\mathit{ls}}} u^{-c}$ và chấp nhận khi
$\Hash(u, \tilde u', \mathit{nonce}) = c$.

\begin{theorem}[Tính đầy đủ]
Với ví trung thực thì $\tilde u' = \tilde u$.
\end{theorem}

\begin{proof}
\[
S^{\hat v} R^{\widehat{\mathit{ls}}} u^{-c}
= S^{\tilde v + c v'} R^{\widetilde{\mathit{ls}} + c\mathit{ls}}
  \bigl(S^{v'} R^{\mathit{ls}}\bigr)^{-c}
= S^{\tilde v} R^{\widetilde{\mathit{ls}}} = \tilde u .
\]
Mọi thành phần nhân với $c$ triệt tiêu.
\end{proof}

\begin{theorem}[Tính chắc chắn]
Ai trả lời đúng cho hai thách thức khác nhau $c_1 \neq c_2$ trên cùng một $\tilde u$ thì thực
sự biết $\mathit{ls}$.
\end{theorem}

\begin{proof}
Chia hai đẳng thức kiểm cho nhau được
$S^{\hat v_1 - \hat v_2} R^{\widehat{\mathit{ls}}_1 - \widehat{\mathit{ls}}_2} = u^{c_1 - c_2}$.
Đặt $\Delta c = c_1 - c_2 \neq 0$, ta rút ra
$\mathit{ls} = (\widehat{\mathit{ls}}_1 - \widehat{\mathit{ls}}_2)/\Delta c$ và
$v' = (\hat v_1 - \hat v_2)/\Delta c$ --- một mở cam kết hợp lệ của $u$.
\end{proof}

\subsection{Ký mù và gỡ mù}

Issuer chọn số nguyên tố $e \in [2^{596}, 2^{596} + 2^{119})$ với $\gcd(e, \varphi) = 1$, lấy
$v''$ dài 2724 bit, tính $d = e^{-1} \bmod \varphi$ rồi đặt
\begin{equation}
Q = Z \cdot \Bigl( u\, S^{v''} \prod_i R_i^{m_i} \Bigr)^{-1}, \qquad A = Q^{d} \bmod n .
\end{equation}

Ví nhận $(A, e, v'')$, đặt $v = v' + v''$ và lưu \g{credential}{credential} $(A, e, v)$.

\begin{theorem}[Đẳng thức chữ ký CL]
\label{thm:cl}
\begin{equation}
Z = A^{e}\, S^{v}\, R^{\mathit{ls}} \prod_i R_i^{m_i} \pmod n .
\end{equation}
\end{theorem}

\begin{proof}
Vì $de \equiv 1 \pmod{\varphi}$ nên $A^{e} = Q^{de} = Q$. Thay $Q$ và
$u = S^{v'} R^{\mathit{ls}}$:
\[
A^{e} S^{v} R^{\mathit{ls}} \prod_i R_i^{m_i}
= Z \Bigl(S^{v'} R^{\mathit{ls}} S^{v''} \prod_i R_i^{m_i}\Bigr)^{-1}
  S^{v' + v''} R^{\mathit{ls}} \prod_i R_i^{m_i} = Z . \qedhere
\]
\end{proof}

Đây là bất biến trung tâm của toàn hệ thống: sai một thuộc tính, sai link secret, hay chữ ký giả
--- đều không cho ra $Z$.

\section{Công thức giữa Wallet và Verifier}

\subsection{Ngẫu nhiên hoá và làm mù}

Ví lấy $r$ dài 2128 bit rồi đặt
\begin{equation}
A' = A\,S^{r}, \qquad v^{*} = v - er, \qquad e^{*} = e - 2^{596}.
\end{equation}

\begin{lemma}[Đẳng thức được bảo toàn]
\label{lem:rerand}
$Z = A'^{\,e} S^{v^{*}} R^{\mathit{ls}} \prod_i R_i^{m_i}$.
\end{lemma}

\begin{proof}
Từ $A = A' S^{-r}$ suy ra $A^{e} = A'^{\,e} S^{-er}$. Thay vào Định lý~\ref{thm:cl}, phần
$S^{-er}$ gộp vào số mũ của $S$ biến $v$ thành $v - er = v^{*}$.
\end{proof}

Gọi $\Hid$ là tập thuộc tính giữ kín, $\Rev$ là tập tiết lộ. Ví lấy $\tilde e$ (456 bit),
$\tilde v$ (3060 bit), $\tilde m_{\mathit{ls}}$ và $\tilde m_i$ (593 bit mỗi số) rồi tính
\begin{equation}
T = A'^{\,\tilde e} S^{\tilde v} R^{\tilde m_{\mathit{ls}}} \prod_{i \in \Hid} R_i^{\tilde m_i},
\qquad c = \Hash(T, A', n_v),
\end{equation}
\[
\hat e = \tilde e + c e^{*}, \quad
\hat v = \tilde v + c v^{*}, \quad
\hat m_{\mathit{ls}} = \tilde m_{\mathit{ls}} + c\,\mathit{ls}, \quad
\hat m_i = \tilde m_i + c m_i \ (i \in \Hid).
\]

\begin{remark}
Mọi số che đều dài hơn số bị che ít nhất 80 bit. Vì $\tilde s$ phân bố đều trên một khoảng độ
dài $2^{\ell_{\tilde s}}$ còn $c\,s$ chỉ là một hằng số trong khoảng $2^{\ell_s + \ell_c}$,
khoảng cách thống kê giữa $\tilde s + c s$ và số ngẫu nhiên thuần tuý nhỏ hơn $2^{-80}$.
\end{remark}

\subsection{Phép kiểm của verifier}

Verifier kiểm biên ($1 < A' < n$, $0 \le \hat e < 2^{457}$, $0 \le c < 2^{256}$, tập thuộc tính
khớp yêu cầu, $\enc$ của giá trị thô khớp $m_i$), rồi tính
\begin{equation}
D = Z \Bigl(\prod_{i \in \Rev} R_i^{m_i}\Bigr)^{-1} A'^{\,-2^{596}}, \qquad
\widehat T = D^{-c} A'^{\,\hat e} S^{\hat v} R^{\hat m_{\mathit{ls}}}
             \prod_{i \in \Hid} R_i^{\hat m_i},
\end{equation}
và chấp nhận khi $\Hash(\widehat T, A', n_v) = c$.

\begin{lemma}
\label{lem:D}
$D = A'^{\,e^{*}} S^{v^{*}} R^{\mathit{ls}} \prod_{i \in \Hid} R_i^{m_i}$.
\end{lemma}

\begin{proof}
Từ Bổ đề~\ref{lem:rerand} và $e = e^{*} + 2^{596}$, chuyển hai thừa số mà verifier tự tính được
sang vế trái là ra đúng biểu thức của $D$.
\end{proof}

\begin{theorem}[Tính đúng đắn]
Với bằng chứng trung thực, $\widehat T = T$.
\end{theorem}

\begin{proof}
\begin{align*}
A'^{\,\hat e} S^{\hat v} R^{\hat m_{\mathit{ls}}} \prod_{i \in \Hid} R_i^{\hat m_i}
&= A'^{\,\tilde e + c e^{*}} S^{\tilde v + c v^{*}}
   R^{\tilde m_{\mathit{ls}} + c \mathit{ls}}
   \prod_{i \in \Hid} R_i^{\tilde m_i + c m_i} \\
&= \underbrace{\Bigl(A'^{\,\tilde e} S^{\tilde v} R^{\tilde m_{\mathit{ls}}}
   \prod R_i^{\tilde m_i}\Bigr)}_{T}
   \cdot
   \underbrace{\Bigl(A'^{\,e^{*}} S^{v^{*}} R^{\mathit{ls}}
   \prod R_i^{m_i}\Bigr)^{c}}_{D^{c} \text{ (Bổ đề \ref{lem:D})}}
= T D^{c}.
\end{align*}
Nhân hai vế với $D^{-c}$ được $\widehat T = T$, nên hàm băm cho ra đúng $c$.
\end{proof}

\section{Công thức giữa Issuer và Verifier}

Giữa hai bên này không có phép tính chung nào, vì không có kênh liên lạc nào. Toàn bộ quan hệ
toán học giữa họ nằm ở chỗ: verifier sử dụng đúng bộ số $(n, S, R, Z, \{R_i\})$ mà issuer đã
công bố, và đẳng thức chữ ký ở Định lý~\ref{thm:cl} chỉ đúng khi bằng chứng thực sự bắt nguồn
từ một credential do issuer đó ký.

Nói cách khác, blockchain đóng vai kênh truyền \emph{một chiều, bất đồng bộ, không cần tin
tưởng} giữa issuer và toàn bộ verifier trên thế giới.

% ======================= THIẾT KẾ BẢO MẬT =============================
\chapter{Thiết kế bảo mật}

\section{Mỗi tài khoản một link secret duy nhất}

Ngay sau khi người dùng tạo ví, hệ thống sinh đúng một \g{linksecret}{link secret} và một
\g{blinding}{blinding factor} cho tài khoản đó. Hai con số này là
\g{bigint}{số nguyên lớn} ngẫu nhiên, không dẫn xuất từ mật khẩu hay bất kỳ dữ liệu cá nhân nào,
nên không thể đoán ngược.

Hệ quả thiết kế: \g{credential}{credential} không phải là một tấm vé vô danh dùng chung. Nó gắn
chặt với đúng một link secret. Một bản sao credential rơi vào tay người khác là một mớ số vô
nghĩa, vì họ không tạo nổi bằng chứng nếu thiếu link secret tương ứng.

\section{Dữ liệu người dùng chỉ nằm trên máy của họ}

Sau bước \g{ekyc}{eKYC} một lần duy nhất, dữ liệu định danh gốc không được gửi đi đâu nữa. Mọi
lần xác minh về sau chỉ trao đổi bằng chứng toán học.

Điều này khác biệt về bản chất so với mô hình ``gửi ảnh giấy tờ'': ở đây không tồn tại bản sao
nào để mà rò rỉ. Bên xác minh có bị tấn công thì kẻ tấn công cũng chỉ thu được những bằng chứng
đã dùng, vốn đã hết hạn vì gắn với \g{nonce}{nonce} cũ, và không chứa dữ liệu gốc.

\section{Bảo vệ bằng passkey ngay trên thiết bị}

Kho lưu trữ của ví được khoá bằng \g{passkey}{passkey} theo chuẩn WebAuthn --- vân tay, khuôn
mặt, hoặc mã PIN của chính thiết bị. Khoá riêng của passkey nằm trong phần cứng bảo mật của máy
và không trích xuất ra được.

Cần nói rõ ranh giới: passkey bảo vệ \emph{quyền truy cập} vào ví. Nó không tham gia vào phép
toán tạo bằng chứng. Hai lớp này độc lập nhau, và việc nhầm lẫn giữa chúng là hiểu sai phổ biến
nhất về hệ thống.

Vì bằng chứng gắn với thiết bị đang giữ link secret, người dùng \textbf{không cần} quay video
sinh trắc hay chụp ảnh khuôn mặt mỗi lần xác minh. Một bằng chứng là đủ để khẳng định ``tôi là
tôi''.

\section{Thiết kế ba bên công khai}

Vai trò của ba bên được định nghĩa công khai và ràng buộc bằng mật mã, không bằng thoả thuận:

\begin{center}
\begin{tabular}{@{}p{2.6cm}p{5.2cm}p{6.4cm}@{}}
\toprule
\textbf{\sffamily Bên} & \textbf{\sffamily Làm được gì} & \textbf{\sffamily Không làm được gì} \\
\midrule
\g{issuer}{Issuer} &
Ký credential cho người đã qua eKYC &
Không thấy link secret; không biết thẻ dùng ở đâu; không giả được thẻ cho người chưa eKYC \\
\addlinespace
\g{holder}{Holder} &
Tạo bằng chứng, chọn trường chia sẻ &
Không sửa được thuộc tính đã ký; không tạo được bằng chứng cho credential mình không giữ \\
\addlinespace
\g{verifier}{Verifier} &
Kiểm bằng chứng, biết chắc thẻ hợp lệ &
Không thấy trường không được mở; không ghép được hai lần trình bày \\
\bottomrule
\end{tabular}
\end{center}

Không bên nào phải tin bên nào. Uy tín giữa các bên đến từ việc mọi ràng buộc đều kiểm chứng
được độc lập, bằng dữ liệu công khai trên \g{blockchain}{blockchain}.

% ======================= THIẾT KẾ BLOCKCHAIN ==========================
\chapter{Thiết kế blockchain}

\section{Issuer chỉ công bố đúng hai thứ}

Issuer ghi lên chain duy nhất \g{pubkey}{khoá công khai} và \g{schema}{schema}. Không gì khác.

\subsection{Khoá công khai gồm những gì}

Chi tiết đầy đủ ở mục~\ref{sec:setup}. Nói cho dễ hiểu: đó là một bộ số cho phép bất kỳ ai
\emph{kiểm} được chữ ký của issuer, nhưng không cho phép ai \emph{tạo} ra chữ ký đó --- giống
như con dấu mà ai cũng nhận ra được nhưng không ai khắc lại được.

\subsection{Schema là gì}

Schema là khung mô tả một loại giấy tờ có những trường nào. Ví dụ với thẻ sinh viên:

\begin{center}
\begin{tabular}{@{}ll@{}}
\toprule
\textbf{\sffamily Trường} & \textbf{\sffamily Ý nghĩa} \\
\midrule
\texttt{mssv} & Mã số sinh viên \\
\texttt{name} & Họ và tên \\
\texttt{class} & Mã lớp \\
\texttt{expiry} & Ngày hết hạn \\
\bottomrule
\end{tabular}
\end{center}

Trong bản demo hiện tại, schema được dùng là mặt trước thẻ căn cước công dân với tám trường: số
thẻ, họ và tên, ngày sinh, giới tính, quốc tịch, quê quán, nơi thường trú, ngày hết hạn.

\section{Verifier dùng schema để soạn yêu cầu}

Vì schema nằm công khai trên chain, verifier biết chính xác loại giấy tờ này có gì, từ đó soạn
được \g{presreq}{yêu cầu trình bày} hợp lệ. Với thẻ sinh viên ở trên, verifier có thể hỏi:

\begin{itemize}[leftmargin=1.5em,itemsep=4pt]
  \item Người này có mã số sinh viên do trường cấp không?
  \item Thẻ sinh viên còn hạn không?
  \item Mã lớp có đúng là lớp đang thi không?
\end{itemize}

Verifier chỉ cần lấy schema từ chain về là đủ để làm việc --- không phải liên hệ, không phải ký
hợp đồng, không phải chờ ai cấp quyền.

\section{Những gì không lên chain}

\begin{center}
\begin{tabular}{@{}p{6.4cm}p{7.4cm}@{}}
\toprule
\textbf{\sffamily Có trên chain} & \textbf{\sffamily Không lên chain} \\
\midrule
Khoá công khai của issuer & Dữ liệu cá nhân của bất kỳ ai \\
Schema của từng loại giấy tờ & Credential đã cấp \\
Danh sách issuer được tin cậy & Link secret, blinding factor \\
& Lịch sử ai xác minh cái gì, ở đâu \\
\bottomrule
\end{tabular}
\end{center}

Lý do cần đến blockchain chỉ có một: khoá công khai phải nằm ở nơi \emph{không ai sửa được}, kể
cả chính issuer. Nếu khoá đặt trên một máy chủ thường, người kiểm soát máy chủ đó có thể âm thầm
thay khoá rồi ký thẻ giả mà không ai phát hiện. Blockchain ở đây đóng vai một \emph{bảng tin
công khai không thể tẩy xoá}, không phải một cơ sở dữ liệu.

% ======================= USE CASE =====================================
\chapter{Các tình huống ứng dụng}

\section{Bhutan NDI --- đã triển khai thực tế}

\g{ndi}{Bhutan NDI} là hệ thống định danh số quốc gia của vương quốc Bhutan, xây trên tiêu chuẩn
\g{anoncreds}{AnonCreds} và đã phục vụ khoảng 200.000 người dân.

Giá trị mang lại:

\begin{itemize}[leftmargin=1.5em,itemsep=5pt]
  \item \textbf{Rút ngắn thời gian thủ tục.} Quy trình mở tài khoản ngân hàng vốn mất nhiều
        ngày chờ đối chiếu hồ sơ nay rút xuống còn vài phút, vì bằng chứng kiểm được tại chỗ.
  \item \textbf{Giảm gánh nặng lưu trữ cho doanh nghiệp.} Ngân hàng và nhà cung cấp dịch vụ
        không còn phải lưu bản sao giấy tờ, kéo theo giảm rủi ro pháp lý và chi phí tuân thủ.
  \item \textbf{Người dân giữ quyền kiểm soát.} Mỗi lần chia sẻ là một lần người dân chủ động
        đồng ý, với đúng những trường họ chọn.
\end{itemize}

\section{Chống thi hộ tại các trường đại học Việt Nam}

Tình trạng gian lận danh tính trong thi cử --- thi hộ, mượn thẻ sinh viên --- diễn ra khá phổ
biến. Các biện pháp hiện nay đều có điểm yếu: kiểm thẻ giấy thì dễ làm giả, đối chiếu ảnh thì
tốn thời gian và sai sót, còn thu thập sinh trắc học thì lại tạo ra một kho dữ liệu nhạy cảm
mới.

\subsection{Cách NX Cred giải quyết}

Trước kỳ thi, nhà trường (đóng vai \g{issuer}{issuer}) cấp cho mỗi sinh viên một
\g{credential}{credential} chứa mã số sinh viên, mã lớp, ngành và ngày hết hạn. Tại phòng thi,
giám thị (đóng vai \g{verifier}{verifier}) tạo một \g{presreq}{yêu cầu} với vài mệnh đề đơn
giản:

\begin{itemize}[leftmargin=1.5em,itemsep=4pt]
  \item Có tên trong danh sách dự thi của môn này không?
  \item Có đúng mã lớp không?
  \item Thẻ sinh viên còn hiệu lực không?
\end{itemize}

Sinh viên quét mã QR, ví tạo bằng chứng, giám thị nhận kết quả trong vài giây.

\subsection{Vì sao cơ chế này chặn được thi hộ}

Điểm then chốt nằm ở \g{linksecret}{link secret}. Muốn nhờ người khác thi hộ, sinh viên phải
giao cả ví lẫn link secret cho người đó. Nhưng link secret không phải một mã số dùng riêng cho
kỳ thi --- nó là thứ mở khoá \emph{toàn bộ} credential trong ví: thẻ ngân hàng, giấy tờ tuỳ
thân, mọi giấy tờ số khác.

Nói cách khác, cái giá của việc nhờ thi hộ là trao toàn bộ danh tính số của mình cho người
khác. Chi phí đó đủ lớn để cơ chế tự bảo vệ chính nó, mà không cần tới camera hay sinh trắc học.

\section{Ngân hàng và dịch vụ tài chính}

Quy trình định danh khách hàng của ngân hàng hiện nay tốn kém ở cả hai phía: khách hàng phải
nộp hồ sơ nhiều lần cho nhiều ngân hàng, còn ngân hàng phải lưu trữ và bảo vệ kho dữ liệu đó
suốt vòng đời quan hệ khách hàng.

Với NX Cred, khách hàng làm \g{ekyc}{eKYC} một lần với cơ quan có thẩm quyền. Sau đó mỗi ngân
hàng chỉ cần một \g{zkp}{bằng chứng} rằng khách hàng có giấy tờ hợp lệ, do đúng cơ quan nhà nước
cấp, và chưa hết hạn. Ngân hàng đạt được mục đích tuân thủ mà không phải trở thành nơi lưu trữ
dữ liệu nhạy cảm --- giảm đồng thời chi phí vận hành và rủi ro pháp lý khi có sự cố.

% ======================= KẾT LUẬN =====================================
\chapter{Kết luận}

NX Cred cho thấy mô hình định danh không tiết lộ là khả thi trên thực tế, không chỉ trên lý
thuyết. Hệ thống đạt được đồng thời ba điều mà mô hình truyền thống buộc phải đánh đổi: bên xác
minh tin được kết quả, người dùng không lộ dữ liệu thừa, và không bên nào theo dõi được người
dùng.

Cụ thể, bản hiện thực này đã chạy được: \g{blindsign}{ký mù} bằng
\g{clsig}{chữ ký CL} thật; \g{sigma}{sigma protocol} khi xin cấp thẻ;
\g{selective}{tiết lộ chọn lọc} theo từng trường; \g{rerand}{ngẫu nhiên hoá} chữ ký tạo
\g{unlink}{tính không liên kết được}; và công bố khoá công khai lên
\g{blockchain}{blockchain} để verifier kiểm độc lập.

\section{Giới hạn hiện tại}
\label{sec:limits}

Tài liệu này sẽ không trung thực nếu bỏ qua những gì hệ thống \emph{chưa} làm được.

\begin{description}[leftmargin=0pt,style=nextline,itemsep=8pt]
  \item[Bằng chứng còn đơn giản]
    Hiện ứng dụng mới chứng minh được một mệnh đề duy nhất: ``credential này đã được issuer
    ký''. Những mệnh đề phong phú hơn --- \g{predicate}{bằng chứng vị từ} kiểu ``đủ 18 tuổi mà
    không lộ ngày sinh'', hay ``thuộc một trong các lớp sau'' --- đòi hỏi bổ sung thêm thuật
    toán tính toán, cụ thể là range proof và set membership proof. Riêng với chứng minh tuổi,
    còn phải đổi cách mã hoá ngày sinh sang dạng số có thứ tự, vì hiện tại ngày sinh bị băm nên
    mất hoàn toàn quan hệ so sánh.

  \item[Chưa có cơ chế thu hồi]
    Issuer chưa \g{revocation}{thu hồi} được credential đã cấp. Hệ quả trực tiếp: mất thiết bị
    đồng nghĩa mất toàn bộ credential, và phải xin cấp lại từ đầu. Nghiêm trọng hơn, một
    credential cấp nhầm hoặc cấp cho người sau đó mất tư cách vẫn tiếp tục có hiệu lực.

  \item[Ví chưa chạy hoàn toàn phía người dùng]
    Trong bản demo, phần mật mã của ví chạy trên máy chủ để rút ngắn thời gian dựng sản phẩm.
    Đúng tinh thần \g{ssi}{định danh tự chủ} thì phần này phải chạy trong trình duyệt, để
    \g{linksecret}{link secret} không bao giờ rời khỏi thiết bị người dùng.
\end{description}

\section{Hướng phát triển tương lai}

\begin{enumerate}[leftmargin=1.6em,itemsep=7pt]
  \item \textbf{Giao diện soạn bằng chứng cho người không chuyên.}
    Mục tiêu là verifier chỉ cần bấm chọn từ danh sách mệnh đề dựng sẵn --- ``đủ 18 tuổi'',
    ``còn hiệu lực'', ``thuộc lớp X'' --- và hệ thống tự sinh yêu cầu tương ứng, thay vì phải
    hiểu cấu trúc \g{schema}{schema} và tự soạn.

  \item \textbf{Bổ sung cơ chế thu hồi.}
    Thêm sổ đăng ký thu hồi (revocation registry) trên chain, để issuer vô hiệu hoá được
    credential và ví chứng minh được ``credential của tôi chưa bị thu hồi'' mà vẫn không lộ
    danh tính.

  \item \textbf{Bằng chứng vị từ.}
    Hiện thực range proof cho các thuộc tính số, mở đường cho nhóm mệnh đề so sánh vốn là nhu
    cầu phổ biến nhất trong thực tế.

  \item \textbf{Chuyển lõi mật mã về trình duyệt.}
    Biên dịch phần tính toán sang WebAssembly để chạy hoàn toàn phía người dùng, hoàn thiện mô
    hình tự chủ danh tính.
\end{enumerate}

% ======================= THAM KHẢO ====================================
\chapter*{Tham khảo}
\addcontentsline{toc}{chapter}{Tham khảo}

\begin{enumerate}[leftmargin=1.6em,itemsep=5pt]
  \item AnonCreds Specification --- \url{https://hyperledger.github.io/anoncreds-spec/}
  \item J. Camenisch, A. Lysyanskaya. \emph{A Signature Scheme with Efficient Protocols}, 2002.
  \item W3C. \emph{Verifiable Credentials Data Model} ---
        \url{https://www.w3.org/TR/vc-data-model-2.0/}
  \item W3C. \emph{Web Authentication (WebAuthn) Level 3} ---
        \url{https://www.w3.org/TR/webauthn-3/}
  \item Bhutan National Digital Identity --- \url{https://www.bhutanndi.com/}
  \item Hyperledger Indy --- \url{https://hyperledger-indy.readthedocs.io/}
  \item A. Fiat, A. Shamir. \emph{How to Prove Yourself: Practical Solutions to Identification
        and Signature Problems}, 1986.
\end{enumerate}

% ======================= PHỤ LỤC ======================================
\appendix

\chapter{Các thư viện sử dụng}

\section{Phía máy chủ --- Python}

\begin{longtable}{@{}p{2.8cm}p{6.4cm}p{4.6cm}@{}}
\toprule
\textbf{\sffamily Thư viện} & \textbf{\sffamily Vai trò} & \textbf{\sffamily Địa chỉ} \\
\midrule
\endhead
gmpy2 & Số học \g{bigint}{số nguyên lớn}: sinh \g{safeprime}{safe prime}, luỹ thừa modulo. Gánh toàn bộ phần nặng của lõi mật mã. & \url{https://pypi.org/project/gmpy2/} \\
\addlinespace
FastAPI & Khung web dựng các endpoint của issuer, ví và verifier. & \url{https://fastapi.tiangolo.com/} \\
\addlinespace
Uvicorn & Máy chủ chạy ứng dụng FastAPI. & \url{https://www.uvicorn.org/} \\
\addlinespace
Pydantic & Kiểm tra và ràng buộc kiểu dữ liệu vào ra của API. & \url{https://pypi.org/project/pydantic/} \\
\addlinespace
py\_webauthn & Xác minh \g{passkey}{passkey} theo chuẩn WebAuthn. & \url{https://pypi.org/project/webauthn/} \\
\addlinespace
google-auth & Kiểm ID-token khi đăng nhập bằng tài khoản Google. & \url{https://pypi.org/project/google-auth/} \\
\addlinespace
web3.py & Đọc khoá công khai và schema từ hợp đồng trên chain. & \url{https://pypi.org/project/web3/} \\
\addlinespace
websockets & Kênh thời gian thực giữa máy tính và điện thoại khi quét giấy tờ. & \url{https://pypi.org/project/websockets/} \\
\bottomrule
\end{longtable}

\section{Phía trình duyệt --- TypeScript}

\begin{longtable}{@{}p{2.8cm}p{6.4cm}p{4.6cm}@{}}
\toprule
\textbf{\sffamily Thư viện} & \textbf{\sffamily Vai trò} & \textbf{\sffamily Địa chỉ} \\
\midrule
\endhead
React & Dựng giao diện ví và cổng verifier. & \url{https://react.dev/} \\
\addlinespace
Vite & Công cụ đóng gói và máy chủ phát triển. & \url{https://vite.dev/} \\
\addlinespace
TypeScript & Ngôn ngữ có kiểu tĩnh, bắt lỗi trước khi chạy. & \url{https://www.typescriptlang.org/} \\
\addlinespace
Tailwind CSS & Hệ thống tạo kiểu giao diện. & \url{https://tailwindcss.com/} \\
\addlinespace
Framer Motion & Hiệu ứng chuyển cảnh giữa các bước. & \url{https://motion.dev/} \\
\addlinespace
jsQR & Đọc mã QR in trên mặt trước giấy tờ. & \url{https://www.npmjs.com/package/jsqr} \\
\addlinespace
Tesseract.js & Đọc chữ trên ảnh giấy tờ, bù các trường mã QR không chứa. & \url{https://tesseract.projectnaptha.com/} \\
\addlinespace
qrcode & Sinh mã QR chuyển phiên từ máy tính sang điện thoại. & \url{https://www.npmjs.com/package/qrcode} \\
\addlinespace
React Router & Điều hướng giữa các màn hình. & \url{https://reactrouter.com/} \\
\bottomrule
\end{longtable}

\chapter{Mã nguồn}

Mã nguồn đầy đủ của ứng dụng được lưu tại kho Git của dự án. Bảng dưới đây đối chiếu từng phần
lý thuyết trong tài liệu với vị trí hiện thực tương ứng.

\begin{center}
\begin{tabular}{@{}p{6.2cm}p{7.4cm}@{}}
\toprule
\textbf{\sffamily Nội dung} & \textbf{\sffamily Vị trí trong mã nguồn} \\
\midrule
Sinh khoá, schema, ký mù & \texttt{issuer/issuer.py} \\
Link secret, cam kết, sigma protocol & \texttt{wallet/wallet.py} \\
Ngẫu nhiên hoá và tạo bằng chứng & \texttt{wallet/wallet.py} \\
Phép kiểm bằng chứng & \texttt{verifier/verifier.py} \\
Mã hoá thuộc tính, hàm băm thách thức & \texttt{common.py} \\
Đọc ghi hợp đồng trên chain & \texttt{chain.py} \\
Xác minh passkey (WebAuthn) & \texttt{passkey.py} \\
Giao diện ví và cổng verifier & \texttt{frontend/src/apps/} \\
\bottomrule
\end{tabular}
\end{center}

\chapter{Video demo và hướng dẫn sử dụng}

{\itshape\color{muted} Phần này sẽ được bổ sung sau.}

\end{document}
